\usepackage[utf8]{inputenc} % This line sets the input encoding to "utf8" for proper handling of special characters
\usepackage[T1]{fontenc} % This line sets the font encoding to "T1"
\usepackage[french]{babel} % This line sets the default language to French
\usepackage{enumitem} % This line imports the "enumitem" package for customizing lists
%\usepackage{multicol} % This line imports the "multicol" package for multi-column layout
\usepackage{underscore} % This line imports the "underscore" package for handling underscores in text
%\usepackage{lipsum} % Import du package lipsum pour générer du texte fictif (lorem ipsum)
\usepackage[autolanguage]{numprint} % This line imports the "numprint" package with the option "autolanguage" for automatic language detection in number formatting
\usepackage{fancyvrb} % This line imports the "fancyvrb" package for customizing verbatim environments
\usepackage{xcolor} % This line imports the "xcolor" package for custom colors
%\usepackage{wallpaper} % This line imports the "wallpaper" package for adding background images
\usepackage{mdframed} % This line imports the "mdframed" package for creating framed boxes
\usepackage{wallpaper}

%%% Line spacing
\usepackage{setspace} % This line imports the "setspace" package for controlling line spacing
%\onehalfspacing % Sets the line spacing to 1.5

%%% Maths
%\usepackage{amsmath} % Import du package amsmath pour les commandes mathématiques avancées
%\usepackage{amsfonts} % Import du package amsfonts pour les polices mathématiques supplémentaires
%\usepackage{amsthm} % Import du package amsthm pour les environnements de théorème et de preuve

%%% Bibliography
\usepackage{biblatex} % This line is commented out and imports the "biblatex" package for bibliography management
\addbibresource{biblio.bib} % This line is commented out and adds the bibliography resource file "fin.bib" for citation management

%%% Graphics
\usepackage{graphicx} % Import du package graphicx pour l'insertion et la manipulation d'images
\usepackage{tikz} % This line imports the "tikz" package for creating graphics and diagrams
\graphicspath{ {assets/bgs/} }

%%% HyperLinks
\usepackage{xurl} % pour casser les URL longs
\usepackage{hyperref} % pour rendre les URL cliquables
\hypersetup{
	colorlinks=true, % Sets the hyperlinks to be colored
	linkcolor=blue, % Sets the color for internal links
	urlcolor=red, % Sets the color for URLs
}

%%% Geometry
%\usepackage{geometry} % Import du package geometry pour la personnalisation de la géométrie de la page
%\geometry{
%	paper=a4paper, % Sets the paper size to A4
%	inner=1cm, % Sets the inner margin to 1cm
%	outer=1cm, % Sets the outer margin to 1cm
	%bindingoffset=.5cm, % Sets the binding offset to 0.5cm (extra space on the left side for binding)
%	top=1.5cm, % Sets the top margin to 1cm
%	bottom=1.5cm, % Sets the bottom margin to 1cm
	%showframe, % Uncomment to show how the type block is set on the page
%}
\usepackage[a4paper, margin=2cm]{geometry}
%%% Lists / Codes
% \usepackage{listings} % This line imports the "listings" package for typesetting source code
% \lstdefinestyle{c}{
% 	language=c, % Sets the language to C for proper syntax highlighting
% 	basicstyle=\ttfamily\small, % Sets the basic font style to monospace and small size
% 	commentstyle=\color{green!40!black}, % Sets the style for comments to green color
% 	keywordstyle=\color{blue}, % Sets the style for keywords to blue color
% 	numberstyle=\tiny\color{gray}, % Sets the style for line numbers to tiny size and gray color
% 	numbers=left, % Positions the line numbers on the left side
% 	stepnumber=1, % Sets the step size for line numbers to 1
% 	numbersep=5pt, % Sets the distance between line numbers and code to 5pt
% 	backgroundcolor=\color{white}, % Sets the background color of the code block to white
% 	frame=single, % Adds a frame around the code block
% 	breaklines=true, % Enables automatic line breaking
% 	breakatwhitespace=true, % Breaks lines at any whitespace
% 	showspaces=false, % Hides spaces with special characters
% 	showstringspaces=false, % Hides spaces within strings
% 	showtabs=false, % Hides tab characters
% 	morekeywords={True, False}, % Adds additional keywords to be highlighted
% 	literate={ê}{{\^e}}1 {é}{{'e}}1 {è}{{e}}1 {à}{{a}}1 {ç}{{\c c}}1 {ï}{{i}}1 {ï}{{i}}1 {î}{{\c i}}1 {ï}{{`i}}1 {ô}{{\^o}}1 {û}{{\^u}}1 {â}{{\^a}}1 {ë}{{"e}}1 {ö}{{"o}}1 {ü}{{"u}}1 {ä}{{"a}}1 {_}{\_}1, % Allows special characters to be displayed correctly
% }

%%% Chapters
% \usepackage{emptypage} % This line imports the "emptypage" package to avoid numbering empty pages
% \usepackage{fancyhdr} % This line imports the "fancyhdr" package for custom page headers and footers
% \usepackage{etoolbox} % This line imports the "etoolbox" package for patching commands
% \patchcmd{\chapter}{\thispagestyle{plain}}{\thispagestyle{fancy}}{}{} % This line patches the command for chapter headings to use the "fancy" page style instead of the "plain" style

%%% Chapters titles
% \usepackage{titlesec} % This line imports the "titlesec" package for customizing section and chapter titles
% \titleformat
% 	{\chapter} % This line defines the "chapter" environment as a chapter
% 	[display] % This line defines the "display" style for the chapter title
% 	{\bfseries\Large\itshape} % This line sets the font style and size for the chapter title
% 	{\begin{center}\vspace{-2cm}Partie \thechapter\end{center} } % This line sets the label for the chapter title
% 	{-0.5cm} % This line sets the spacing between the chapter title and the label
% 	% {vspace{0.5ex}} % This line sets the vertical spacing between the chapter title and the label
% 	{
% 		\rule{\textwidth}{0.5pt} % Adds a horizontal rule with a thickness of 0.5pt
% 		%\vspace{0.25ex} % Adds vertical space of 0.25ex
% 		\centering % Centers the chapter title
% 	} % before-code
% 	[
% 		\vspace{-0.5ex} % Removes vertical space of 0.5ex
% 		\rule{\textwidth}{0.05pt} % Adds a thin horizontal rule with a thickness of 0.05pt
% 	] % after-code

%%% Sections & Subsections numbering
\renewcommand{\thesection}{\arabic{section}}
\renewcommand{\thesubsection}{\arabic{section}.\arabic{subsection}}

%%% Theorems
%\usepackage[most]{tcolorbox} % This line is commented out and imports the "tcolorbox" package for creating colored boxes
%\newtcbtheorem{theo}% % This line is commented out and defines a new theorem-like environment called "theo"
%{Test for convergence}{}{theorem} % The environment name is "theo", the heading is "Test for convergence", and the counter is labeled as "theorem"

\date{} % Sets the date to empty